\section{Waste}
\label{sec:waste}

\subsection{Scope and inventory mapping}

This block targets the \textbf{area-source} component of CAMS GNFR~J
(\emph{Waste}). In the CAMS-REG-ANT inventory, GNFR~J is reported as a
single aggregated waste sector, while the fine-scale geography of waste
emissions is known to differ substantially between landfill-related
activity, wastewater handling, and diffuse residual waste processes.
The methodological problem is therefore not the construction of final
downscaled emissions at this stage, but the construction of
pollutant-specific \emph{within-cell spatial weights} that can later be
combined with the corresponding CAMS cell totals.

This distinction is essential. The workflow described here produces a
multi-band raster of normalised fine-grid weights, one band per
pollutant, such that the fine-pixel shares sum to unity within each
CAMS GNFR~J area cell. It does not by itself alter the CAMS emission
mass, nor does it claim that CAMS already resolves the internal waste
subsectors needed for spatial allocation.

\subsection{Conceptual design}

GNFR~J is internally heterogeneous. Landfill-type activity, composting,
anaerobic digestion, domestic and industrial wastewater handling, open
waste burning, and other diffuse waste processes do not follow the same
fine-scale spatial pattern. A single land-cover mask would therefore
force all waste-related pollutants to share the same geography despite
their different process origins.

The upgraded design separates the problem into three linked elements:
\begin{itemize}
  \item a set of three operational proxy families representing solid
        waste treatment, wastewater, and residual diffuse waste;
  \item pollutant-specific country-level family shares derived from
        CEIP/NRD waste subsector totals;
  \item a final composite proxy that recombines the common family
        surfaces with the pollutant-specific national shares and is then
        normalised within each CAMS parent cell.
\end{itemize}

This produces a transparent two-stage weighting system. First, the
ancillary geospatial datasets describe where different waste-related
activities are most plausibly located. Second, the CEIP/NRD national
subsector totals determine how strongly each pollutant should follow
each proxy family in each country. Pollutant bands may therefore differ
spatially even though CAMS provides only one aggregate GNFR~J area layer.
The operational family split and the supporting waste proxy datasets are
summarised in Appendix~\ref{app:waste_downscaling},
Tables~\ref{tab:waste_families} and \ref{tab:waste_proxy_inputs}.

\subsection{Operational proxy families and pollutant shares}

The detailed CEIP/NRD waste subsectors available for this block include
biological treatment of waste, waste incineration, open burning,
wastewater handling, and residual waste categories. For the area-proxy
workflow, these subsectors are aggregated into three operational
families:
\begin{itemize}
  \item \textbf{Solid waste treatment:} 5A, 5B1, 5B2;
  \item \textbf{Wastewater:} 5D1, 5D2, 5D3;
  \item \textbf{Residual diffuse waste:} 5C2, 5E.
\end{itemize}

Incineration-like subsectors (\textbf{all 5C1$\ast$ codes}) are excluded
from the default area-weighting workflow. This is a deliberate design
choice intended to avoid double counting with point-source treatment of
waste incineration facilities, which are conceptually closer to
stack-like infrastructure than to diffuse area activity.

Let $c$ denote country and $p$ pollutant. The national CEIP/NRD family
totals are written as
\begin{equation}
  E_{\mathrm{solid},c,p} =
  E_{5A,c,p} + E_{5B1,c,p} + E_{5B2,c,p},
\end{equation}
\begin{equation}
  E_{\mathrm{ww},c,p} =
  E_{5D1,c,p} + E_{5D2,c,p} + E_{5D3,c,p},
\end{equation}
\begin{equation}
  E_{\mathrm{res},c,p} =
  E_{5C2,c,p} + E_{5E,c,p}.
\end{equation}

The corresponding pollutant-specific family shares are
\begin{equation}
  w_{\mathrm{solid},c,p} =
  \frac{E_{\mathrm{solid},c,p}}
       {E_{\mathrm{solid},c,p} + E_{\mathrm{ww},c,p} + E_{\mathrm{res},c,p}},
\end{equation}
\begin{equation}
  w_{\mathrm{ww},c,p} =
  \frac{E_{\mathrm{ww},c,p}}
       {E_{\mathrm{solid},c,p} + E_{\mathrm{ww},c,p} + E_{\mathrm{res},c,p}},
\end{equation}
\begin{equation}
  w_{\mathrm{res},c,p} =
  \frac{E_{\mathrm{res},c,p}}
       {E_{\mathrm{solid},c,p} + E_{\mathrm{ww},c,p} + E_{\mathrm{res},c,p}}.
\end{equation}

By construction,
\begin{equation}
  w_{\mathrm{solid},c,p} +
  w_{\mathrm{ww},c,p} +
  w_{\mathrm{res},c,p} = 1.
\end{equation}

If the denominator is zero for a country--pollutant pair, the
implementation uses a transparent fallback ladder: first the mean family
mixture from other pollutants available for the same country, then a
European mean across countries with valid data, and finally an equal
split across the three families. These fallbacks should be interpreted
as regularisation devices that keep the weighting system operational
when subsector-resolved waste reporting is incomplete.

\subsection{Family-specific spatial proxies}

Let $i$ denote a fine pixel on the reference grid. The method constructs
three non-negative family surfaces
$P_{\mathrm{solid},c}(i)$, $P_{\mathrm{ww},c}(i)$, and
$P_{\mathrm{res},c}(i)$ from harmonised ancillary datasets, reprojected
to a common fine grid and transformed to unitless relevance fields.
These surfaces do not represent emissions directly; they represent the
relative spatial plausibility of each family.

\paragraph{Solid waste treatment.}
The solid-waste proxy is anchored primarily in CORINE Land Cover class
132 (\emph{dump sites}), which is the most direct harmonised European
land-cover indicator of landfill-type activity. Because composting and
anaerobic digestion are not consistently captured by one dedicated
pan-European land-cover class, the method adds CORINE class 121
(\emph{industrial or commercial units}) as a secondary support layer.
OpenStreetMap waste-related features, where available, are used only as
geometric refinement of this support and not as a standalone estimate of
activity intensity. This hierarchy reflects the stronger semantic link
between class~132 and solid waste disposal, while acknowledging the weak
harmonised EU-wide evidence base for composting and anaerobic digestion.

\paragraph{Wastewater.}
The wastewater proxy combines wastewater infrastructure and settlement
support. UWWTD agglomerations describe the broad serviced population
footprint, while UWWTD treatment plants provide infrastructure anchors
for treatment-related activity. These layers are complemented by GHSL
population and Copernicus imperviousness as distributed support for
built-up wastewater generation. An optional industrial-land term may be
added to reflect the spatial relevance of industrial wastewater
handling. The resulting proxy therefore links domestic wastewater mainly
to populated serviced areas, while retaining explicit support from
wastewater infrastructure.

\paragraph{Residual diffuse waste.}
The residual diffuse proxy is designed to remain spatially diffuse
rather than infrastructure-dominated. It combines GHSL population with
GHSL settlement classes emphasising rural or low-density settlement
patterns, together with imperviousness as a supplementary settlement
indicator where appropriate. This reflects the intended interpretation
of residual diffuse waste as a combination of general settlement-related
waste activity and open burning processes that are more plausibly
associated with rural or peri-urban populations than with large
specialised facilities.

\subsection{Composite proxy and within-cell normalisation}

For pollutant $p$, country $c$, and fine pixel $i$, the raw composite
waste proxy is
\begin{equation}
  P_{J,c,p}(i) =
  w_{\mathrm{solid},c,p} \, P_{\mathrm{solid},c}(i)
  + w_{\mathrm{ww},c,p} \, P_{\mathrm{ww},c}(i)
  + w_{\mathrm{res},c,p} \, P_{\mathrm{res},c}(i).
\end{equation}

This equation expresses the central logic of the workflow. The
fine-scale geography comes from the three common family surfaces, while
the pollutant-specific and country-specific differences come from the
CEIP/NRD-derived weights. A pollutant whose national waste budget is
dominated by wastewater will inherit more of the wastewater geography;
a pollutant dominated by landfill-type or residual waste processes will
inherit more of those respective families.

Let $m$ denote a CAMS GNFR~J area cell and let $i \in m$ denote fine
pixels whose centres fall inside that CAMS cell. The final within-cell
weight is then
\begin{equation}
  W_{J,c,p}(i) =
  \frac{P_{J,c,p}(i)}
       {\displaystyle\sum_{k \in m} P_{J,c,p}(k)}.
\end{equation}

By construction,
\begin{equation}
  \sum_{i \in m} W_{J,c,p}(i) = 1
  \qquad \forall m,p.
\end{equation}

If the denominator is zero for a given CAMS cell and pollutant, the
algorithm falls back to a uniform distribution over the valid fine-grid
support of that CAMS cell. This ensures that a CAMS cell retains a
complete weight field even when the fine-scale proxy evidence is absent
locally.

The final product is therefore a set of pollutant-specific
\emph{normalised weighting rasters}. These layers are intended for a
later downscaling step in which they may be combined with the
corresponding CAMS cell totals. They are not themselves the downscaled
emissions.

\subsection{Inputs and software}

The main inputs are summarised in
Table~\ref{tab:waste_inputs}. The underlying datasets themselves are
described in Appendix~\ref{app:datasets}, including the CEIP/NRD waste
workbook (Appendix~\ref{app:waste_ceip}), the UWWTD layers
(Appendix~\ref{app:waste_uwwtd}), the GHSL population and settlement
products (Appendix~\ref{app:waste_ghsl}), Copernicus imperviousness
(Appendix~\ref{app:waste_imperviousness}), and the optional
OpenStreetMap waste-feature extract (Appendix~\ref{app:waste_osm}).

\begin{table}[htbp]
  \centering
  \caption{Main inputs for CAMS GNFR~J waste weighting layers.}
  \label{tab:waste_inputs}
  \begin{tabular}{p{3.4cm}p{9.3cm}}
    \hline
    Ingredient & Role \\
    \hline
    CAMS-REG-ANT NetCDF & GNFR~J area-cell geometry, pollutant list, and parent-cell support for within-cell normalisation \\
    CEIP/NRD waste emissions & Country- and pollutant-specific subsector totals used to derive $w_{\mathrm{solid}}$, $w_{\mathrm{ww}}$, and $w_{\mathrm{res}}$ \\
    CORINE CLC 2018 & Reference grid; class~132 for dump sites and class~121 for industrial/commercial support \\
    OpenStreetMap waste features & Optional geometric refinement of solid-waste support \\
    UWWTD agglomerations & Spatial support for serviced wastewater population \\
    UWWTD treatment plants & Infrastructure anchors for wastewater activity \\
    GHSL population / settlement classes & Diffuse settlement support for wastewater and residual waste families \\
    Copernicus imperviousness & Built-surface support for wastewater and residual waste proxies \\
    NUTS geometry & Country masks and reference-grid domain construction \\
    \hline
  \end{tabular}
\end{table}

The implementation is designed as a modular Python workflow using
\texttt{rasterio}, \texttt{geopandas}, \texttt{xarray},
\texttt{numpy}, \texttt{pandas}, and related GIS utilities. The output
is written as a multi-band GeoTIFF of pollutant-specific weights,
accompanied by diagnostic tables for the CEIP family shares and for
CAMS cells requiring zero-proxy fallback.

\subsection{Interpretation and limitations}

The three family surfaces should be interpreted as first-order spatial
likelihood layers rather than as direct measurements of waste activity.
They identify plausible environments of occurrence for the main waste
process groups represented inside GNFR~J, but they do not reconstruct
facility throughput, wastewater loads, or waste quantities in physical
units at pixel level.

Several limitations follow directly from that role.

First, CAMS provides only an aggregate GNFR~J area layer and does not
separate the waste sector into internally resolved area subcategories.
The decomposition into three proxy families is therefore an external
analytical construction rather than an inventory-native distinction.

Second, harmonised European proxy evidence is uneven across the waste
subsectors. Landfill-type activity is relatively well represented by
CORINE class~132, but composting and anaerobic digestion do not have an
equally direct EU-wide land-cover equivalent. Their spatial placement
therefore relies partly on broader industrial/commercial support and, if
available, auxiliary OSM geometry.

Third, waste incineration subsectors are excluded by default from the
area-weighting methodology. This improves consistency with a separate
point-source treatment of incineration facilities, but it also means
that the area proxy is intentionally not a full representation of all
reported CEIP waste processes.

Despite these limitations, the GNFR~J workflow represents a substantial
improvement over a single undifferentiated waste proxy. It preserves the
country- and pollutant-specific subsector structure visible in CEIP/NRD
reporting, allows landfill, wastewater, and diffuse waste processes to
follow different fine-scale geographies, and maintains exact
within-CAMS-cell normalisation for the weighting raster ultimately used
in the downscaling pipeline.
